\documentclass[12pt]{article}
\usepackage[T1]{fontenc}
\usepackage[utf8]{inputenc}
\usepackage{lmodern}
\usepackage{textcomp}
\usepackage{enumitem}
\usepackage{geometry}
\geometry{margin=1in}
\begin{document}
\begin{center}
Answer Key - Sociology - Graduate Level Assessment 13 \
\small (Answers are ordered exactly as in the question paper.)
\end{center}

\section*{Multiple Choice}
\begin{enumerate}[label=\arabic*.]
\item \textbf{Answer:} C
\\ \textit{Options:} A) marketization and privatization; B) individual freedom, choice and responsibility; C) generous financial benefits for single parents, students and the unemployed; D) the rolling back of the welfare state
\\ \textit{Note:} Thatcher's government prioritized neoliberal economic policies that emphasized reducing state welfare provisions and promoting individual responsibility, which led to a decrease in financial benefits for single parents, students, and the unemployed rather than striving for generosity in these areas.
\item \textbf{Answer:} C
\\ \textit{Options:} A) reading a textual document quickly to understand the gist of it; B) gathering a sample from whoever you can see in a public place; C) scanning a table to identify general patterns and significant figures; D) keeping your eyes on research participants at all times
\\ \textit{Note:} 'Eyeballing' refers to the practice of visually scanning data, such as tables or graphs, to quickly discern overarching trends and noteworthy data points, which is essential for sociologists to comprehend and analyze complex information effectively.
\item \textbf{Answer:} C
\\ \textit{Options:} A) setting and streaming pupils by ability; B) parental choice of school; C) supporting LEAs that appeared to be failing; D) state regulation and control of the curriculum
\\ \textit{Note:} The New Labour government, under Tony Blair, focused on reforming education through initiatives that emphasized standards, choice, and accountability, often sidelining local education authorities (LEAs) that were perceived as failing, rather than actively supporting them.
\item \textbf{Answer:} A
\\ \textit{Options:} A) decreased rates of prostitution and sex tourism; B) developing countries can depend on it as a crucial source of income; C) the exploitation of cheap, unregulated labour in poor countries; D) we have become more aware of 'other' societies and ways of living
\\ \textit{Note:} The correct answer is based on the understanding that global tourism often leads to increased rates of prostitution and sex tourism as destinations cater to the demands of tourists, making it unlikely for these activities to decrease as a consequence of tourism.
\item \textbf{Answer:} D
\\ \textit{Options:} A) it is committed on a larger, often global, scale, and is well organized; B) it is associated with political conflict between states and their citizens; C) it can have far-reaching effects upon international relations; D) all of the above
\\ \textit{Note:} The correct answer "all of the above" is justified because terrorism encompasses political motives, collective action, and aims to instill fear on a broader societal level, which distinguishes it fundamentally from the individualistic and localized crime patterns emphasized by the Chicago School of sociology.
\end{enumerate}

\section*{True / False}
\begin{enumerate}[label=\arabic*.]
\setcounter{enumi}{5}
\item \textbf{Answer:} False
\\ \textit{Options:} A) True; B) False
\\ \textit{Note:} While impersonal rules are important in a bureaucracy, Weber's ideal type emphasizes rationality and hierarchical organization rather than solely focusing on impersonal rules as a fundamental component.
\item \textbf{Answer:} False
\\ \textit{Options:} A) True; B) False
\\ \textit{Note:} The nuclear family specifically refers to a household unit consisting of two parents and their children, rather than an extended network of relatives.
\item \textbf{Answer:} False
\\ \textit{Options:} A) True; B) False
\\ \textit{Note:} The 'new racism' refers to a form of racial discrimination that is often subtle and embedded in cultural and social norms, rather than an anti-fascist movement opposing nationalist politics, which suggests a misunderstanding of its nature and implications in contemporary society.
\item \textbf{Answer:} True
\\ \textit{Options:} A) True; B) False
\\ \textit{Note:} Telework involves employees performing their job duties remotely, typically from home, using digital communication tools to connect with their employer, thus making the statement true.
\item \textbf{Answer:} False
\\ \textit{Options:} A) True; B) False
\\ \textit{Note:} This statement is false because systemic inequalities, such as gender biases, discrimination, and unequal access to resources and opportunities, continue to create significant barriers for women in both employment and political representation.
\end{enumerate}

\section*{Long Form Response}
\begin{enumerate}[label=\arabic*.]
\setcounter{enumi}{10}
\item \textbf{Answer:} The Conservative government of 1979 implemented several strategies to diminish the influence of the labour movement, most notably through legislation that made all strike action illegal. This drastic measure aimed to undermine the bargaining power of trade unions and limit workers' ability to collectively assert their rights. By criminalizing strike actions, the government effectively curtailed the capacity of workers to protest against unfair labor practices or unsafe working conditions, thereby shifting the balance of power significantly in favor of employers. The implications of such strategies were profound, as they not only suppressed workers' rights but also fostered a climate of fear and repression within the labor movement, leading to a significant erosion of collective bargaining and solidarity among workers. Ultimately, these policies reflect broader ideological shifts towards neoliberalism, prioritizing market interests over social welfare and workers' rights.
\\ \textit{Note:} correct because it identifies the key strategy of criminalizing strike actions as a means to weaken trade unions and highlights the consequent erosion of workers' rights and ability to address grievances.
\item \textbf{Answer:} Social stratification refers to the hierarchical arrangement of individuals within society, often based on factors such as wealth, income, education, and social status. This structure creates distinct layers, which can significantly influence individual mobility- defined as the ability to move between these layers. In societies with rigid stratification systems, such as caste or class-based societies, the mobility of individuals is often limited, leading to persistent inequalities and social tensions. Conversely, more fluid societies may offer greater opportunities for upward mobility, yet still face challenges related to systemic barriers that affect marginalized groups. Overall, social stratification not only shapes the distribution of resources and power but also impacts social cohesion and the overall dynamics of societal change.
\\ \textit{Note:} The answer correctly identifies social stratification as a hierarchical system that impacts individual mobility and highlights the consequences of rigid stratification on social inequality.
\end{enumerate}

\vfill
\noindent\textit{End of Answer Key}
\end{document}